%============================================================
% CH1 · PATTERN 4 — SIMULATION AND INDEX ARITHMETIC: THEORY
%============================================================
\subsection{Pattern Theory}
\label{ssec:ch1-p4-theory}

\begin{patternbox}{Pattern 4 at a Glance}
Many array problems need no clever idea at all: the statement \textbf{describes a process}, and the solution performs it step by step on the array. The skill is to do it \emph{exactly} --- the right order of updates, the right index formula, the right behaviour at the borders --- and, when memory is tight, to store an old and a new value in the same cell.

\begin{center}
\begin{tikzpicture}[x=1cm, y=1cm]
    \dparray{d}{1.2}{digits}{1/acell, 2/acell, 9/aread, 9/aread}
    \dparray{e}{0}{$+1$}{1/akeep, 3/awrite, 0/awrite, 0/awrite}
    \foreach \i/\j in {3/2, 2/1}{ \draw[dep] ([yshift=2mm]d\i.north) to[bend right=50] node[lbl, above]{carry} ([yshift=2mm]d\j.north); }
    \begin{scope}[xshift=5.6cm]
        \dparray{p}{1.2}{$i$}{0/acell, 1/acell, 2/acell, 3/acell, 4/acell, 5/acell}
        \draw[dep] (p1.south) to[bend right=35] (p2.south);
        \draw[dep] (p1.south) to[bend right=40] (p3.south);
        \node[note, anchor=north] at (2.7,0.45) {$i \to 2i,\ 2i{+}1$ \quad $i \to (i+1) \bmod n$ \quad $(r,c) \to rC + c$};
        \node[note, anchor=north] at (2.7,-0.05) {cell $= \text{old} + n \cdot \text{new}$};
    \end{scope}
\end{tikzpicture}
\end{center}

\begin{recurrencebox}
\centering
\textbf{simulate}: state $\leftarrow$ rule(state), one step at a time, in the order the statement dictates\\[3pt]
\textbf{carry}: $s = a_i + b_i + c$,\; digit $= s \bmod 10$,\; $c = \lfloor s/10 \rfloor$,\; from the least significant end\\[3pt]
\textbf{encode}: $a[i] \leftarrow a[i] + n\cdot \text{new}_i$ \;(both $< n$),\; later old $= a[i] \bmod n$, new $= \lfloor a[i]/n \rfloor$
\end{recurrencebox}
Cost: (number of steps) $\times$ (cost of one step). If that is too much, look for a period or a closed formula.
\end{patternbox}

%------------------------------------------------------------
\subsubsection*{From Pattern 3 to Pattern 4}

Pattern~3 moved values to new positions given by \emph{one} formula, $i \to (i+k) \bmod n$, and showed that the order of the moves matters: writing in place too early destroys values that are still needed. Pattern~4 generalises both halves of that lesson. The rule can be any process the statement describes --- adding one to a number, passing an object around a circle, following $i \to a[i]$ --- and the index formula can be any arithmetic on indices: $2i$, $i/2$, $i \bmod n$, $r \cdot C + c$. The ``don't destroy unread values'' problem gets a new, general fix: \textbf{encode} old and new value in one cell (Lesson~\ref{ssec:ch1-inplace} announced it).

%------------------------------------------------------------
\subsubsection*{When to recognise it --- and when not}

Use Pattern 4 when
\begin{itemize}[leftmargin=*]
    \item the statement says ``repeat the following'', ``every second'', ``in order from $0$ to $n-1$ do \ldots'' and the number of steps times the cost per step fits the budget;
    \item a number is given as an \textbf{array of digits} (too long for \texttt{long long}): add, increment, compare, multiply by a small value;
    \item positions are computed: pairs $(2i, 2i+1)$, a circle ($\bmod n$), a grid flattened to 1D, a permutation ($i \to a[i]$);
    \item the array must be rewritten \textbf{in place} while every cell still needs its old value (``$\bigO{1}$ extra memory'' follow-ups);
    \item the input is small and the answer is ``what does the process produce?'' --- simulate first, think later.
\end{itemize}

\textbf{Do NOT use it} blindly when the number of steps is huge ($10^{9}$ seconds of passing a pillow): find the period or the formula (E4.7). Do not simulate a process that repeatedly \emph{inserts into the middle} or \emph{deletes} from a large array --- each step is $\bigO{n}$; that calls for other structures (queues, \sqhbook{2}{Implementing Stacks and Queues}; linked arrays, Chapter~\ref{ch:linked_arrays}). When steps interact only through sums, prefix sums or difference arrays (Chapter~\ref{ch:prefix_sums}) replace the simulation.

%------------------------------------------------------------
\subsubsection*{Tool 1 --- simulation done right}

A simulation is a loop whose body applies the rule once. Three questions decide whether it is correct:
\begin{enumerate}[leftmargin=*]
    \item \step{What is the state?} Sometimes the whole array, often much less (a position and a direction; a running counter that replaces ``decrease all later cells by 1'', E4.5).
    \item \step{In which order are the cells updated?} If the new value of a cell depends on old values of other cells, updating in the wrong order mixes old and new values. Either iterate in an order where the needed values are still old (right to left for a shift to the right), keep a copy, or encode (Tool~4).
    \item \step{What happens at the borders?} First and last index, an empty array, a carry out of the last digit, a pillow that reaches the end of the line.
\end{enumerate}

\begin{diagrambox}
\begin{tikzpicture}[x=1cm, y=1cm]
    \begin{scope}
        \node[font=\small\bfseries, text=deepblue, anchor=west] at (-0.6,1.15) {shift right, left to right \no};
        \dparray{w}{0.3}{}{7/acell, 7/awrite, 7/awrite, 7/awrite}
        \foreach \i/\j in {0/1, 1/2, 2/3}{ \draw[dep, warningred] (w\i.south) to[bend right=40] (w\j.south); }
        \node[note] at (1.8,-0.75) {$a[i] = a[i-1]$ for $i = 1, 2, 3$ on $[7,8,9,5]$};
    \end{scope}
    \begin{scope}[xshift=6.6cm]
        \node[font=\small\bfseries, text=deepblue, anchor=west] at (-0.6,1.15) {shift right, right to left \yes};
        \dparray{r}{0.3}{}{7/acell, 7/awrite, 8/awrite, 9/awrite}
        \foreach \i/\j in {2/3, 1/2, 0/1}{ \draw[dep, successgreen] (r\i.south) to[bend right=40] (r\j.south); }
        \node[note] at (1.8,-0.75) {$i = 3, 2, 1$: each source is read before it is overwritten};
    \end{scope}
\end{tikzpicture}
\figcap{The same update rule in two orders. Left to right copies the first value everywhere; right to left is correct.}
\end{diagrambox}

%------------------------------------------------------------
\subsubsection*{Tool 2 --- digit arrays and carries}

A number too long for any integer type is stored as an array of digits. LeetCode stores it \textbf{most significant digit first} ($123 \to [1,2,3]$); many implementations prefer \textbf{least significant first}, because then digit $i$ has weight $10^i$ and the number can grow at the end with \texttt{push\_back}. Addition is the school algorithm: from the least significant digit, add the digits and the carry, write $s \bmod 10$, carry $\lfloor s / 10 \rfloor$. With two digits and an incoming carry of at most 1, $s \le 9 + 9 + 1 = 19$, so the carry is always 0 or 1.

\begin{diagrambox}
\begin{tikzpicture}[x=1cm, y=1cm]
    \dparray{x}{2.0}{$a$}{9/acell, 9/acell, 8/acell}
    \dparray{y}{1.0}{$b$}{{}/emptycell, 4/acell, 7/acell}
    \draw[thick] (-0.3,0.45) -- (2.9,0.45);
    \begin{scope}[xshift=-0.9cm]
        \dparray{z}{-0.1}{}{1/awrite, 0/awrite, 4/awrite, 5/awrite}
    \end{scope}
    \node[rowlabel] at (-1.05,-0.1) {sum};
    \foreach \x in {-0.45, 0.45, 1.35}{ \node[font=\scriptsize\bfseries, text=accentorange] at (\x,2.6) {1}; }
    \node[font=\scriptsize\bfseries, text=accentorange] at (2.25,2.6) {0};
    \node[note, anchor=west, text width=6.4cm, align=left] at (3.6,1.0)
        {$8 + 7 = 15$: write 5, carry 1; \; $9 + 4 + 1 = 14$: write 4, carry 1; \; $9 + 0 + 1 = 10$: write 0, carry 1; \; the final carry becomes a new leading digit: $998 + 47 = 1045$.\\[3pt]
        Orange: the carry \emph{into} each column from its right.};
\end{tikzpicture}
\figcap{Addition of digit arrays: the result can be one digit longer than the longer input --- the only border case.}
\end{diagrambox}

\begin{invariantbox}[Invariant of the carry loop]
After processing the $t$ least significant columns, the written digits $r_0, \dots, r_{t-1}$ and the carry $c$ satisfy
\[
\sum_{i<t} r_i\,10^{i} \;+\; c\cdot 10^{t} \;=\; \sum_{i<t} (a_i + b_i)\,10^{i}.
\]
\textbf{Proof.} True for $t = 0$ ($c = 0$). Column $t$ computes $s = a_t + b_t + c$ and writes $r_t = s \bmod 10$, $c' = \lfloor s/10 \rfloor$, so $r_t\,10^t + c'\,10^{t+1} = s\,10^t = (a_t + b_t)10^t + c\,10^t$: the old carry term is replaced exactly. When both numbers are exhausted and $c = 0$, the written digits are the sum. \hfill$\blacksquare$
\end{invariantbox}

Incrementing (Problem~\ref{ssec:ch1-plus-one}) is the special case $b = 1$: the carry can only continue through digits equal to 9, so the loop can stop at the first digit that is not 9.

%------------------------------------------------------------
\subsubsection*{Tool 3 --- index arithmetic}

Most of the difficulty of ``easy'' simulations is the index formula. The common ones:

\begin{table}[H]
\centering\small
\renewcommand{\arraystretch}{1.2}
\begin{tabularx}{\linewidth}{@{}>{\raggedright\arraybackslash}p{3.3cm} >{\raggedright\arraybackslash}p{4.3cm} X@{}}
\toprule
\textbf{Map} & \textbf{Meaning} & \textbf{Typical use} \\
\midrule
$i \to a[i]$ & follow a pointer stored in the cell & permutations, ``next'' arrays (Problem~\ref{ssec:ch1-build-permutation}) \\
$i \to 2i,\ 2i+1$; \ $i \to \lfloor i/2 \rfloor$ & split into / merge from pairs & interleaving two halves, blocks of 2 (E4.2), heaps \\
$i \to (i \pm 1) \bmod n$ & neighbours on a circle & circular arrays, Josephus-type walks \\
$(r, c) \leftrightarrow rC + c$ & row-major 2D $\leftrightarrow$ 1D & reshaping, grids as arrays (E4.1, E4.3) \\
$i \to n-1-i$ & mirror & palindromes, reversal \\
$t \to t \bmod 2(n-1)$ & back-and-forth walk has period $2(n-1)$ & bouncing objects (E4.7) \\
\bottomrule
\end{tabularx}
\end{table}

\begin{diagrambox}
\begin{tikzpicture}[x=1cm, y=1cm]
    \begin{scope}
        \node[font=\small\bfseries, text=deepblue, anchor=west] at (-0.6,1.0) {grid $3 \times 4$};
        \foreach \r in {0,1,2}{ \foreach \c in {0,...,3}{
            \pgfmathtruncatemacro{\id}{\r*4+\c}
            \ifnum\id=9 \node[aread] (g\r\c) at (\c*0.9, -\r*0.75) {\id}; \else \node[cell] (g\r\c) at (\c*0.9, -\r*0.75) {\id}; \fi
        }}
    \end{scope}
    \node[note, anchor=west, text width=6.5cm, align=left] at (4.0,-0.75) {cell $(2,1)$ holds id $2 \cdot 4 + 1 = 9$;\\ back: row $\lfloor 9/4 \rfloor = 2$, column $9 \bmod 4 = 1$};
    \draw[flow] (1.35,-2.0) -- node[lbl, right]{row-major} (1.35,-2.7);
    \begin{scope}[xshift=-0.6cm, yshift=-3.3cm]
        \dparray{f}{0}{}{0/acell, 1/acell, 2/acell, 3/acell, 4/acell, 5/acell, 6/acell, 7/acell, 8/acell, 9/aread, 10/acell, 11/acell}
        \arange{f0}{f3}{row 0}{memoorange}
        \arange{f4}{f7}{row 1}{memoorange}
        \arange{f8}{f11}{row 2}{memoorange}
    \end{scope}
\end{tikzpicture}
\figcap{A grid with $C$ columns stored as one array: $(r,c) \mapsto rC + c$ and back with $/$ and $\bmod$ by $C$ (the number of \emph{columns}, not rows).}
\end{diagrambox}

\textbf{Negative indices.} In C++ \texttt{-7 \% 5 == -2}. For a position on a circle use \texttt{((i \% n) + n) \% n}. And remember that ``number of cells from $l$ to $r$'' is $r - l + 1$ for a closed range but $r - l$ for a half-open one (Lesson~\ref{ssec:ch1-indices}).

%------------------------------------------------------------
\subsubsection*{Tool 4 --- two values in one cell}

Suppose every cell must become a new value that depends on \emph{old} values of other cells, and we may not allocate a second array. If all values (old and new) lie in $[0, n)$, store
\[
a[i] \;=\; \text{old}_i + n \cdot \text{new}_i ,
\]
like a two-digit number in base $n$. While the cells are being rewritten, the old value of \emph{any} cell is still available as $a[j] \bmod n$, whether $j$ has been processed or not; at the end a second pass divides by $n$.

\begin{diagrambox}
\begin{tikzpicture}[x=1cm, y=1cm]
    \node[anscell, minimum width=22mm, minimum height=9mm, font=\large\bfseries] (big) at (0,0) {13};
    \node[note] at (0,-0.85) {one cell, $n = 6$};
    \node[basecell, minimum width=14mm] (lo) at (4.2,0.55) {$13 \bmod 6 = 1$};
    \node[compcell, minimum width=14mm] (hi) at (4.2,-0.55) {$\lfloor 13 / 6 \rfloor = 2$};
    \draw[flow] (big.east) -- (lo.west); \draw[flow] (big.east) -- (hi.west);
    \node[font=\scriptsize\bfseries, anchor=west] at (5.6,0.55) {old value (low digit)};
    \node[font=\scriptsize\bfseries, anchor=west] at (5.6,-0.55) {new value (high digit)};
    \node[note, anchor=west, text width=3.6cm, align=left] at (8.7,0) {$1 + 6 \cdot 2 = 13$; the maximum is $(n-1) + n(n-1) = n^2 - 1$};
\end{tikzpicture}
\figcap{Base-$n$ encoding of an (old, new) pair. Problem~\ref{ssec:ch1-build-permutation} uses exactly this cell.}
\end{diagrambox}

\textbf{Conditions.} (1) Both values lie in $[0, n)$ --- otherwise use any $M$ larger than every value, $a[i] = \text{old} + M \cdot \text{new}$. (2) $M^2$ fits in the type: $n = 10^5$ gives $10^{10}$, which needs \texttt{long long}. (3) When writing cell $i$, read the old value of the source cell \emph{through} \texttt{\% M}, because the source may already have been encoded. Variants of the same idea: bit packing (\texttt{a[i] |= new << 10} when values are below $2^{10}$), and the \textbf{sign} as a one-bit flag when all values are positive (Chapter~\ref{ch:index_tricks}).

%------------------------------------------------------------
\subsubsection*{Tool 5 --- boundary handling}

\begin{itemize}[leftmargin=*]
    \item \textbf{Guard every computed index:} \texttt{if (0 <= j \&\& j < n)} before \texttt{a[j]}; out-of-range reads in C++ do not crash reliably, they return garbage.
    \item \textbf{Sentinels:} an extra cell at one or both ends (a $0$ in front of a digit array, $-\infty$ around an array of heights) removes special cases from the loop.
    \item \textbf{Growth at the front:} a carry out of the most significant digit needs \texttt{insert(begin(), 1)} ($\bigO{n}$, done once) --- or store digits least significant first and \texttt{push\_back}.
    \item \textbf{Clamping:} ``never below 0'' is \texttt{max(0, x)}, applied at the moment of the update, not at the end.
\end{itemize}

\begin{diagrambox}
\begin{tikzpicture}[x=1cm, y=1cm]
    \dparray{s}{0}{}{0/emptycell, 9/acell, 9/acell, 9/acell}
    \dparray{t}{-1.2}{}{1/awrite, 0/awrite, 0/awrite, 0/awrite}
    \node[aptrlbl, text=darkgray, above] at ([yshift=1mm]s0.north) {sentinel};
    \draw[flow] (4.0,-0.6) -- node[lbl, above]{$+1$} (5.2,-0.6);
    \node[note, anchor=west, text width=5.6cm, align=left] at (5.4,-0.6) {with a leading $0$ cell the carry out of $999$ lands in an existing cell; without it, the array must grow at the front (Problem~\ref{ssec:ch1-plus-one})};
\end{tikzpicture}
\figcap{A sentinel turns the only border case of incrementing into an ordinary step (at the price of a possible leading zero to strip).}
\end{diagrambox}

%------------------------------------------------------------
\subsubsection*{Variations map}

\begin{diagrambox}
\begin{tikzpicture}[
    var/.style={rectangle, rounded corners=2pt, draw=patternborder, fill=patternbg,
                font=\scriptsize, align=center, text width=3.25cm, minimum height=1.25cm, inner sep=3pt}]
    \node[var, fill=answerfill, draw=accentorange, text width=5.2cm, font=\scriptsize\bfseries] (root) at (0,0)
        {Simulation \& index arithmetic\\do exactly what the statement says};
    \node[var] (v1) at (-5.4,-2) {\textbf{digit arrays}\\carry, increment, add (Problem~\ref{ssec:ch1-plus-one}, E4.4, E4.6)};
    \node[var] (v2) at (-1.8,-2) {\textbf{follow $i \to a[i]$}\\apply a permutation (Problem~\ref{ssec:ch1-build-permutation})};
    \node[var] (v3) at (1.8,-2)  {\textbf{2D $\leftrightarrow$ 1D}\\$rC + c$, reshape (E4.1, E4.3)};
    \node[var] (v4) at (5.4,-2)  {\textbf{blocks and pairs}\\$\lfloor i/2 \rfloor$, $2i$, $2i+1$ (E4.2)};
    \node[var] (v5) at (-3.6,-3.6) {\textbf{lazy state}\\one counter replaces ``update all later cells'' (E4.5)};
    \node[var] (v6) at (0,-3.6)    {\textbf{period / formula}\\too many steps: $t \bmod 2(n-1)$ (E4.7)};
    \node[var] (v7) at (3.6,-3.6)  {\textbf{in place}\\old $+\,n\cdot$new in one cell (Problem~\ref{ssec:ch1-build-permutation})};
    \node[draw=patternborder, dashed, rounded corners=4pt, fit=(v1)(v4)(v5)(v7), inner sep=5pt] (grp) {};
    \draw[flow] (root.south) -- (grp.north);
\end{tikzpicture}
\figcap{The boxes differ in what the state is and how indices are computed; the loop is always ``apply the rule''.}
\end{diagrambox}

%------------------------------------------------------------
\subsubsection*{Code templates}

\begin{lstlisting}[style=cpp]
#include <bits/stdc++.h>
using namespace std;

// sum of two numbers given as digit arrays, most significant digit first
vector<int> addDigits(const vector<int>& a, const vector<int>& b) {
    vector<int> res;                                  // built least significant first
    int i = (int)a.size() - 1, j = (int)b.size() - 1, carry = 0;
    while (i >= 0 || j >= 0 || carry) {
        int s = carry;
        if (i >= 0) s += a[i--];
        if (j >= 0) s += b[j--];
        res.push_back(s % 10);
        carry = s / 10;
    }
    reverse(res.begin(), res.end());
    return res;
}

// in place: a[i] <- a[p[i]] for a permutation p of 0..n-1, values of a in [0, n)
void applyInPlace(vector<int>& a, const vector<int>& p) {
    int n = a.size();
    for (int i = 0; i < n; i++) a[i] += n * (a[p[i]] % n);   // new value as the high "digit"
    for (int i = 0; i < n; i++) a[i] /= n;
}

// walk around a circle of n seats, `steps` moves to the right (steps may be negative)
int circular(int start, long long steps, int n) {
    return (int)(((start + steps) % n + n) % n);
}

int main() {
    vector<int> s = addDigits({9, 9, 8}, {4, 7});
    for (int d : s) cout << d;
    cout << "\n";                                     // 998 + 47 = 1045
    vector<int> a = {2, 3, 0, 1}, p = {3, 0, 1, 2};
    applyInPlace(a, p);
    for (int x : a) cout << x << " ";
    cout << "\n";                                     // a[p[i]]: 1 2 3 0
    cout << circular(2, -7, 5) << " " << circular(4, 1000000000000LL, 5) << "\n";  // 0 4
    int C = 4;                                        // 2D <-> 1D, row-major, C columns
    int id = 2 * C + 1;                               // cell (2,1)
    cout << id << " -> (" << id / C << "," << id % C << ")\n";   // 9 -> (2,1)
}
\end{lstlisting}

\begin{insightbox}[Check the templates by hand]
$a = [2,3,0,1]$, $p = [3,0,1,2]$: the new array is $[a_3, a_0, a_1, a_2] = [1,2,3,0]$. During the first pass cell 1 needs the old $a_0 = 2$, but cell 0 already holds $2 + 4 \cdot 1 = 6$; reading $6 \bmod 4 = 2$ recovers it. On the circle: $2 - 7 = -5 \equiv 0 \pmod 5$, and $4 + 10^{12} \equiv 4$ because $10^{12}$ is a multiple of 5. The four printed lines come from running the program.
\end{insightbox}

%------------------------------------------------------------
\subsubsection*{Common mistakes}

\begin{warningbox}[Mistake 1 --- updating in an order that reads new values]
\texttt{for (i = 1; i < n; i++) a[i] = a[i-1];} copies $a[0]$ everywhere. Find the order in which every source is still old, or use a copy, or encode.
\end{warningbox}

\begin{warningbox}[Mistake 2 --- converting a digit array to an integer]
100 digits do not fit in \texttt{long long} (19 digits) or even \texttt{unsigned \_\_int128} (38 digits). Work digit by digit.
\end{warningbox}

\begin{warningbox}[Mistake 3 --- reading an encoded cell without \texttt{\% n}]
\begin{lstlisting}[style=snippet]
a[i] += n * a[p[i]];          // WRONG: a[p[i]] may already be old + n*new
a[i] += n * (a[p[i]] % n);    // RIGHT: always the old value
\end{lstlisting}
\end{warningbox}

\begin{warningbox}[Mistake 4 --- overflow of the encoding]
$\text{old} + n\cdot\text{new} \le n^2 - 1$. For $n = 1000$ that is about $10^6$ (fine in \texttt{int}); for $n = 10^5$ it is about $10^{10}$ and needs \texttt{long long}.
\end{warningbox}

\begin{warningbox}[Mistake 5 --- rows and columns mixed up]
For an $R \times C$ grid the 1D index is $rC + c$ and the inverse uses $/C$ and $\%C$. Using $R$ works only for square grids, so square test cases hide the bug.
\end{warningbox}

\begin{checklistbox}[``I'm stuck'' checklist for Pattern 4]
\begin{itemize}[leftmargin=*]
    \item[$\square$] What exactly is the state, and what does one step change? Can I write the step as a 3-line function?
    \item[$\square$] Steps $\times$ cost per step: does it fit? If not, is there a period, a formula, or a counter that replaces an inner loop?
    \item[$\square$] Does a cell's new value need old values of other cells? Which update order keeps them old? Do I need a copy or an encoding?
    \item[$\square$] Is my index formula right on a $2 \times 3$ (non-square, tiny) example?
    \item[$\square$] Borders: first/last index, a carry out of the top digit, negative \texttt{\%}, an empty result.
\end{itemize}
\end{checklistbox}
